% ============================================================
\section{Iteración 7 --- DRV-07: localización de extremo a extremo}\label{sec:it07}
% ============================================================

% ------------------------------------------------------------
\subsection{Driver y alcance}

\begin{tabla}{|L{3.0cm}|Y|}
\fichacab
\parte{Driver}{DRV-07 --- Localización de extremo a extremo.}
\parte{Enunciado}{Todo texto visible, fecha, número y formato se puede traducir y adaptar por región sin tocar código, y los textos viajan en Unicode desde la pantalla hasta las columnas de SQL Server y de vuelta.}
\parte{Por qué le toca este turno}{Empata en 10 con DRV-06 y va después porque agrupa un solo ASR de prioridad alta, ASR-13, que nace de un objetivo de negocio y no de un requisito impuesto con stakeholders de veto. Llega aquí, y no más tarde, porque las tres convenciones que impone atraviesan la interfaz, el protocolo y el esquema, y las tres piezas ya existen: agregarlas después obligaría a revisarlas todas.}
\parte{Qué decide esta iteración}{Cómo se agrega un idioma nuevo sin tocar código y sin rehacer pantallas, y cómo se garantiza que el texto no se dañe en ningún tramo del recorrido.}
\parte{Decisiones ya tomadas}{DEC-01 a DEC-03 y lo decidido en las seis iteraciones anteriores. Pesan CD-17, que ya fijó que el texto viaja en Unicode dentro de mensajes con esquema, y CD-34, que puso el catálogo de palabras prohibidas por idioma fuera del código.}
\parte{Qué no decide}{El contenido de las traducciones ni qué idiomas se publican, que son decisiones de STK-04 y del equipo de localización. Tampoco decide el diseño gráfico, aunque sí le imponga condiciones.}
\end{tabla}

% ------------------------------------------------------------
\subsection{Paso 1. Revisión de entradas}

\subsubsection*{ASR que agrupa el driver}

Del conjunto de la sección~\ref{sec:asr}, este es el ASR que DRV-07 agrupa y el papel que cumple en esta iteración.

\begin{tabla}{|L{1.4cm}|Y|L{1.7cm}|L{3.6cm}|}
\hline
\filacab \cab{ASR} & \cab{Qué debe cumplir la arquitectura} & \cab{Prioridad} & \cab{Papel en esta iteración} \\ \hline
\endhead
ASR-13 & Todo texto y formato se adapta a otro idioma sin tocar código, y el texto viaja en Unicode desde la pantalla hasta la base de datos. & Alta & Objetivo. Su medida es ESC-11. \\ \hline
\end{tabla}

\subsubsection*{Objetivos de diseño}

\begin{tabla}{|L{1.5cm}|L{5.0cm}|Y|}
\hline
\filacab \cab{ID} & \cab{Objetivo} & \cab{Origen y qué exige aquí} \\ \hline
\endhead
OBJ-21 & El juego se distribuye internacionalmente. &
Objetivo declarado del proyecto, con REQ-06. Es el que vuelve obligatorio el cambio de idioma. \\ \hline
OBJ-22 & El equipo de localización trabaja sin depender de los programadores. &
CRN-16, con STK-11. Una traducción que exige tocar código no se entrega dentro de un incremento semanal. \\ \hline
OBJ-11 & Un niño de 8 años entiende la interfaz. &
REQ-07 y CRN-01, con STK-01 y STK-09. Compite con lo anterior: lo más claro para un niño es el ícono con texto dentro (TEN-05). \\ \hline
OBJ-23 & Un jugador puede llamarse como quiera y escribir en su idioma. &
CRN-16 y CON-05, con STK-11. Afecta al protocolo y al esquema, no solo a la interfaz. \\ \hline
\end{tabla}

\subsubsection*{Requisitos funcionales primarios}

\begin{tabla}{|L{1.3cm}|L{4.2cm}|Y|}
\hline
\filacab \cab{CU} & \cab{Caso de uso} & \cab{Qué exige de la localización} \\ \hline
\endhead
CU-12 & Editar el perfil y las preferencias &
Guarda el idioma preferido de la cuenta, que es lo que decide en qué idioma se le habla al jugador en todas las pantallas. \\ \hline
CU-13 & Cambiar el nickname &
El nombre lo escribe el jugador y puede llevar caracteres de cualquier sistema de escritura; después lo ve su rival, que juega en otro idioma. \\ \hline
CU-28 & Enviar un mensaje de chat &
El texto viaja entre jugadores de idiomas distintos y se contrasta contra el catálogo del idioma del canal (RN-01). \\ \hline
CU-21, CU-20 & Colocar un muro; mover el peón &
Los avisos de jugada inválida (FA-02) son texto que cambia de longitud al traducirse y que un niño tiene que entender de un vistazo. \\ \hline
CU-25, CU-36, CU-40 & Finalizar la partida; historial; historial de monedas &
Muestran fechas, números, monedas y contadores, que se escriben distinto en cada región. \\ \hline
\end{tabla}

\subsubsection*{Escenarios de atributos de calidad}

\begin{tabla}{|L{1.7cm}|L{3.3cm}|Y|}
\hline
\filacab \cab{Escenario} & \cab{Atributo} & \cab{Medida} \\ \hline
\endhead
ESC-11 & Flexibilidad: adaptabilidad &
Un idioma con otro sistema de escritura queda listo en menos de una semana con 0 líneas de código modificadas; el 100~\% de los textos visibles se puede traducir y 0 están incrustados en imágenes; con textos un 40~\% más largos, 0 pantallas muestran texto cortado o encimado; y un nombre en japonés hace el viaje completo, cliente, red, SQL Server y pantalla del rival, sin ningún carácter dañado. \\ \hline
ESC-02 & Usabilidad &
Condición. El aviso de jugada inválida tiene que seguir entendiéndose en menos de 5 segundos y con una frase corta, en cualquier idioma. \\ \hline
ESC-04 & Protección del menor &
Condición. El filtro trabaja por idioma, así que un idioma nuevo sin catálogo dejaría el chat sin protección. \\ \hline
\end{tabla}

\subsubsection*{Restricciones}

\begin{tabla}{|L{1.9cm}|L{4.3cm}|Y|}
\hline
\filacab \cab{ID} & \cab{Restricción} & \cab{Cómo limita esta iteración} \\ \hline
\endhead
CON-05 & Los textos deben almacenarse en Unicode & Afecta a los tipos de columna, a la codificación del protocolo y a los archivos de recursos: hay que aplicarla de extremo a extremo. \\ \hline
REQ-06 & El juego debe permitir cambiar de idioma & Prohíbe las cadenas dentro del código y los anchos fijos en la interfaz. \\ \hline
CON-04 & El almacenamiento es SQL Server & El esquema tiene que admitir Unicode en cada columna de texto, incluidos nombres y mensajes. \\ \hline
REQ-07 & Comprensible para niños de 8 años & Empuja hacia los íconos, que es lo que choca con la traducción (TEN-05). \\ \hline
\end{tabla}

\subsubsection*{Decisiones de proyecto ya tomadas}

\begin{tabla}{|L{1.5cm}|L{4.3cm}|Y|}
\hline
\filacab \cab{ID} & \cab{Decisión} & \cab{Qué impone a esta iteración} \\ \hline
\endhead
CD-17 & Serialización con esquema fijo por tipo de mensaje & El texto ya viaja en Unicode dentro de los mensajes; queda comprobar el resto del recorrido. \\ \hline
CD-21 & El motor explica el rechazo & El motivo llega del servidor como un dato, no como una frase, para poder traducirse en el cliente. \\ \hline
CD-34 & Catálogo de palabras prohibidas por idioma & Cada idioma nuevo trae su catálogo además de sus textos. \\ \hline
CD-32 & Acceso a datos aislado tras repositorios & El cambio de tipos de columna a Unicode queda confinado a una capa. \\ \hline
\end{tabla}

\subsubsection*{Concerns}

\begin{tabla}{|L{1.5cm}|Y|}
\hline
\filacab \cab{ID} & \cab{Concern y qué aporta a la iteración} \\ \hline
\endhead
CRN-16 & \emph{Necesito poder traducirlo todo sin tocar el código: textos, fechas, números y formatos. Lo que esté escrito dentro del programa o incrustado en una imagen, yo no lo puedo cambiar.} \\ \hline
CRN-15 & El mismo botón en otro idioma puede ocupar la mitad o el doble; maquetar con anchos fijos descuadra la interfaz al traducir. \\ \hline
CRN-01 & El jugador quiere entender al momento por qué no puede hacer una jugada, lo que empuja hacia lo visual. \\ \hline
\end{tabla}

\subsubsection*{Preguntas abiertas que condicionan la iteración}

\begin{tabla}{|L{1.3cm}|Y|}
\hline
\filacab \cab{ID} & \cab{Cómo condiciona} \\ \hline
\endhead
Q-07 & El nombre definitivo y la identidad visual. La marca es un texto más entre los que deben poder sustituirse sin tocar código. \\ \hline
Q-10 & Si el juego se publica en una plataforma de distribución, que puede exigir idiomas o formatos concretos por región. \\ \hline
\end{tabla}

% ------------------------------------------------------------
\subsection{Paso 2. Objetivo de la iteración y selección de entradas}

\subsubsection*{Priorización de las entradas}

\begin{tabla}{|L{3.2cm}|L{1.3cm}|Y|L{2.2cm}|}
\hline
\filacab \cab{Entrada} & \cab{Peso} & \cab{Por qué pesa lo que pesa} & \cab{Decisión} \\ \hline
\endhead
ESC-11 & Alto & Es la medida del objetivo y reúne las cuatro comprobaciones: tiempo, cobertura, longitud y recorrido completo del texto. & Se atiende \\ \hline
CRN-16, CRN-15 & Alto & Son el origen de la medida y describen las dos formas de fallar: texto que no se puede cambiar y maquetas que se rompen. & Se atiende \\ \hline
CON-05, REQ-06 & Alto & Imponen Unicode de extremo a extremo y el cambio de idioma como requisito. & Se atiende \\ \hline
CU-12, CU-13 & Alto & Uno decide en qué idioma se le habla al jugador; el otro mete texto escrito por el usuario en el recorrido completo. & Se atiende \\ \hline
CON-04 & Alto & Si una columna no admite Unicode, todo lo demás no sirve: el texto se daña al guardarse. & Se atiende \\ \hline
CU-28, CU-21 & Medio & Aportan el texto que viaja entre jugadores y los avisos que hay que traducir sin perder claridad. & Se atiende \\ \hline
ESC-02, REQ-07 & Medio & Condición: la solución no puede volver la interfaz menos clara para un niño (TEN-05). & Se atiende como límite \\ \hline
ESC-04, CD-34 & Medio & Cada idioma nuevo necesita su catálogo, o el chat queda sin protección en ese idioma. & Se atiende como obligación \\ \hline
CU-25, CU-36, CU-40 & Medio & Aportan fechas, números y monedas, que no se traducen sino que se formatean. & Se atiende \\ \hline
Q-07, Q-10 & Bajo & No cambian la estructura; la marca entra como un texto más. & Se atiende como condición \\ \hline
\end{tabla}

\subsubsection*{Objetivo de la iteración}

\begin{tabla}{|L{3.0cm}|Y|}
\fichacab
\parte{Objetivo}{Que se pueda agregar un idioma nuevo, con otro sistema de escritura, sin tocar código ni rehacer pantallas.}
\parte{ASR que satisface}{ASR-13.}
\parte{Driver que cubre}{DRV-07.}
\parte{Medida}{La de ESC-11: el idioma queda listo en menos de una semana con 0 líneas de código modificadas; el 100~\% de los textos visibles se puede traducir y 0 están incrustados en imágenes; con textos un 40~\% más largos, 0 pantallas muestran texto cortado o encimado; y un nombre escrito en japonés recorre cliente, red, SQL Server y pantalla del rival sin ningún carácter dañado.}
\parte{Límite que respeta}{ESC-02 y ESC-04. La interfaz tiene que seguir siendo comprensible para un niño de 8 años, y el chat no puede quedar sin filtro en el idioma nuevo.}
\parte{Incremento que produce}{Una persona ajena al código agrega un idioma y juega una partida completa en él, y un recorrido con pseudolocalización, textos alargados con caracteres especiales, muestra que ninguna pantalla se rompe. La partida se juega además entre dos clientes con idiomas distintos y nombres en sistemas de escritura distintos.}
\parte{Criterio de éxito}{La iteración termina cuando el idioma se agrega sin cambios en el repositorio, el recorrido no encuentra textos cortados y el nombre en japonés se ve igual en las dos pantallas.}
\end{tabla}

% ------------------------------------------------------------
\subsection{Paso 3. Elemento a descomponer}

% ------------------------------------------------------------
\Needspace*{0.4\textheight}
\subsubsection*{EL-07 --- Catálogo de textos y formatos}

\begin{tabla}{|L{2.6cm}|Y|}
\fichacab
\parte{Elemento}{El catálogo de textos y formatos: la pieza que resuelve, para cada idioma y región, qué texto se muestra y cómo se escriben fechas, números y monedas.}
\parte{Qué abarca}{Las cadenas visibles de todas las pantallas, los avisos de error y de jugada inválida, los nombres de estados y formas de término, el formato de fechas, números y monedas, la elección del idioma según la preferencia de la cuenta (CU-12) y la relación entre cada idioma y su catálogo de palabras prohibidas. No abarca el arte ni el sonido, aunque sí les imponga la condición de no llevar texto dentro.}
\parte{Por qué se elige}{Porque la medida del objetivo depende de él: que un idioma se agregue sin tocar código es exactamente la propiedad de este elemento. Se elige frente a descomponer la interfaz del cliente, que es donde el texto se ve pero no donde se decide cuál se muestra, y frente al servicio de datos, que guarda texto pero no lo elige. La parte del recorrido que toca al protocolo y al esquema ya está resuelta por CD-17 y CD-32, y aquí solo se comprueba.}
\parte{Qué falta decidir sobre él}{Dónde viven las cadenas y con qué identificador se piden. Quién decide el idioma de una pantalla, la cuenta o el sistema operativo del equipo. Qué se hace con los textos que el servidor produce, como los motivos de rechazo. Y cómo se evita que una maqueta válida en español se rompa en un idioma más largo.}
\parte{Evidencia en los casos de uso}{CU-12 guarda el idioma preferido de la cuenta, pero ningún caso de uso dice de dónde salen las cadenas que las pantallas muestran. CU-13 permite un nickname escrito por el jugador, que después ve un rival con otro idioma. CU-28 contrasta el texto contra el catálogo del idioma del canal (RN-01), lo que supone que cada idioma tiene el suyo. CU-21 muestra en FA-02 un aviso junto al cursor con el motivo del rechazo: es texto corto, traducible y visto por un niño.}
\parte{Trazabilidad}{DRV-07. CD-17, CD-21, CD-32, CD-34. ESC-11, con ESC-02 y ESC-04 como límites. CU-12, CU-13, CU-20, CU-21, CU-25, CU-28, CU-36, CU-40. CON-05, REQ-06, CON-04, REQ-07. REQ-06. CRN-15, CRN-16, CRN-01. TEN-05. STK-11, STK-08, STK-09, STK-01.}
\end{tabla}

% ------------------------------------------------------------
\subsection{Paso 4. Conceptos de diseño}

% ------------------------------------------------------------
\Needspace*{0.3\textheight}
\subsubsection*{CD-39 --- Cadenas externas identificadas por clave}

\begin{tabla}{|L{2.6cm}|Y|}
\fichacab
\parte{Estilo}{Patrón de internacionalización.}
\parte{Descripción}{Ningún texto visible está escrito en el código: cada uno vive en un catálogo por idioma y se pide por una clave. Atiende ESC-11 en sus dos primeras cifras, el idioma en menos de una semana y el 100~\% de los textos traducibles. Se elige frente a escribir las cadenas donde se usan, que obliga a buscar por todo el código en cada traducción y garantiza que alguna se quede sin traducir.}
\parte{Efectos}{El equipo de localización trabaja sobre archivos, sin abrir el proyecto ni recompilar, que es lo que CRN-16 pide. Obliga a una convención de claves y a una comprobación que detecte claves sin traducir antes de publicar un idioma. Una clave ausente tiene que mostrar algo legible y no romper la pantalla.}
\parte{Trazabilidad}{DRV-07, DRV-08. CD-19. ESC-11. CU-12, CU-20, CU-21, CU-25. REQ-06. CRN-16. Q-07. STK-11, STK-08.}
\end{tabla}

% ------------------------------------------------------------
\Needspace*{0.3\textheight}
\subsubsection*{CD-40 --- El servidor envía códigos, el cliente los traduce}

\begin{tabla}{|L{2.6cm}|Y|}
\fichacab
\parte{Estilo}{Patrón de protocolo aplicado a la presentación.}
\parte{Descripción}{Lo que el servidor manda no son frases sino códigos con sus datos: el motivo de un rechazo, la forma de término de una partida o el estado de una sanción viajan como valores, y el cliente los convierte en texto con el catálogo del idioma del jugador. Atiende ESC-11 en la parte de los textos que nacen en el servidor, que de otro modo quedarían fuera del alcance de la traducción. Se elige frente a que el servidor arme las frases, que lo obligaría a conocer el idioma de cada jugador y volvería a meter texto en el código del servidor.}
\parte{Efectos}{Dos jugadores de una misma partida pueden verla en idiomas distintos. El motivo de rechazo que CD-21 entrega encaja aquí de forma natural, porque ya es un dato. Obliga a mantener sincronizados los códigos del servidor y las claves del catálogo, lo que se comprueba junto con las claves sin traducir.}
\parte{Trazabilidad}{DRV-07, DRV-02. CD-17, CD-21, CD-39. ESC-11, ESC-02. CU-21 (FA-02), CU-25, CU-28, CU-44. REQ-06, CON-05. CRN-16, CRN-01. STK-11, STK-09, STK-05.}
\end{tabla}

% ------------------------------------------------------------
\Needspace*{0.3\textheight}
\subsubsection*{CD-41 --- Unicode comprobado en todo el recorrido}

\begin{tabla}{|L{2.6cm}|Y|}
\fichacab
\parte{Estilo}{Convención transversal, con táctica de capacidad de prueba.}
\parte{Descripción}{El texto se maneja en Unicode en los cuatro tramos, interfaz, protocolo, repositorios y columnas de SQL Server, y una prueba recorre el camino completo con un nombre y un mensaje en un sistema de escritura distinto para comprobar que ningún tramo lo daña. Atiende la última cifra de ESC-11. Se elige frente a confiar en que cada tramo lo haga bien por su cuenta, que es como se descubren estos problemas tarde: basta una columna con el tipo equivocado para perder el texto al guardarlo.}
\parte{Efectos}{Fija el tipo de las columnas de texto del esquema y la codificación del protocolo, que CD-17 ya adelantó. La prueba de recorrido entra en la regresión semanal de CD-22, así que un cambio que rompa el manejo de Unicode se detecta en el incremento en que ocurre.}
\parte{Trazabilidad}{DRV-07, DRV-02, DRV-04. CD-17, CD-22, CD-32. ESC-11, ESC-17. CU-13, CU-28, CU-12. CON-05, CON-04. CRN-16. STK-11, STK-13, STK-10.}
\end{tabla}

% ------------------------------------------------------------
\Needspace*{0.3\textheight}
\subsubsection*{CD-42 --- Formato por región para fechas, números y monedas}

\begin{tabla}{|L{2.6cm}|Y|}
\fichacab
\parte{Estilo}{Patrón de internacionalización.}
\parte{Descripción}{Las fechas, los números y las monedas no se arman con texto: se guardan como valores y se formatean al mostrarlos, según la región del jugador. Atiende la parte de ESC-11 que habla de formatos, que CRN-16 menciona junto a los textos. Se elige frente a escribir el formato a mano en cada pantalla, que produce fechas que un jugador de otra región lee al revés.}
\parte{Efectos}{El historial de CU-36, el de monedas de CU-40 y el fin de partida de CU-25 muestran cifras correctas en cada región sin trabajo extra. Obliga a guardar las fechas en un mismo huso y a convertirlas al mostrarlas, decisión que también afecta a los relojes de partida y a los plazos de sanción.}
\parte{Trazabilidad}{DRV-07, DRV-04. CD-31, CD-39. ESC-11. CU-25, CU-36, CU-40, CU-44. CON-05, REQ-06. REQ-06. CRN-16. STK-11, STK-13.}
\end{tabla}

% ------------------------------------------------------------
\Needspace*{0.3\textheight}
\subsubsection*{CD-43 --- Íconos sin texto y contenedores que crecen}

\begin{tabla}{|L{2.6cm}|Y|}
\fichacab
\parte{Estilo}{Táctica de adaptabilidad aplicada a la interfaz.}
\parte{Descripción}{Ninguna imagen lleva texto dentro, y las pantallas se maquetan con contenedores que crecen con su contenido en lugar de anchos fijos. Atiende la tercera cifra de ESC-11, 0 pantallas rotas con textos un 40~\% más largos, y la cuarta parte de la primera, 0 textos incrustados en imágenes. Se elige frente al ícono con la palabra dibujada dentro, que es lo más claro para un niño y lo único que el equipo de localización no puede cambiar: es la resolución de TEN-05.}
\parte{Efectos}{El arte se vuelve reutilizable entre idiomas, y la marca que Q-07 aún no fija queda separada del resto. El costo es que los íconos sin texto pueden resultar ambiguos, así que se prueban con niños junto con ESC-02. La maquetación flexible es más trabajo de diseño por adelantado.}
\parte{Trazabilidad}{DRV-07, DRV-08. CD-39. ESC-11, ESC-02. CU-20, CU-21 (pasos 3 y 5). REQ-06, REQ-07, CON-15. CRN-15, CRN-01. TEN-05. Q-07. STK-08, STK-09, STK-11, STK-01.}
\end{tabla}

% ------------------------------------------------------------
\Needspace*{0.3\textheight}
\subsubsection*{CD-44 --- Pseudolocalización en la regresión semanal}

\begin{tabla}{|L{2.6cm}|Y|}
\fichacab
\parte{Estilo}{Táctica de capacidad de prueba.}
\parte{Descripción}{Existe un idioma de prueba generado a partir del español, con los textos un 40~\% más largos y con caracteres de otros sistemas de escritura, y el recorrido por las pantallas con ese idioma forma parte de la regresión. Atiende ESC-11 al convertir sus cifras en una comprobación repetible, en lugar de esperar a que llegue un traductor real. Se elige frente a comprobar solo cuando se agrega un idioma de verdad, que descubre los problemas cuando ya hay decenas de pantallas hechas.}
\parte{Efectos}{Un texto escrito en el código o una maqueta con ancho fijo se detectan en el incremento en que se introducen, que es cuando arreglarlos es barato (CON-14). Añade una comprobación más al cierre semanal, junto a la regresión de CD-22 y la medición de CD-27.}
\parte{Trazabilidad}{DRV-07, DRV-08. CD-22, CD-39, CD-43. ESC-11, ESC-17. CU-12, CU-20, CU-21. CON-14, REQ-06. CRN-15, CRN-16, CRN-08. STK-11, STK-10, STK-08.}
\end{tabla}
